% Part: model-theory
% Chapter: basics
% Section: dlo

\documentclass[../../../include/open-logic-section]{subfiles}

\begin{document}

\olfileid{mod}{bas}{dlo}
\section{الترتيبات الخطية الكثيفة}

\begin{defn}
  \emph{الترتيب الخطي الكثيف بلا طرفين} هو !!a{structure}~$\Struct{M}$
  للـ!!{language} التي تحتوي !!{predicate} واحدًا ذا موضعين~$<$ يحقق
  الجمل الآتية:
  \begin{enumerate}
  \item $\lforall[{\OLId{latin-ordinary}{x}}][\lnot {\OLId{latin-ordinary}{x}} < {\OLId{latin-ordinary}{x}}]$؛
  \item $\lforall[{\OLId{latin-ordinary}{x}}][\lforall[{\OLId{latin-ordinary}{y}}][\lforall[{\OLId{latin-ordinary}{z}}][({\OLId{latin-ordinary}{x}} < {\OLId{latin-ordinary}{y}} \lif ({\OLId{latin-ordinary}{y}} < {\OLId{latin-ordinary}{z}} \lif {\OLId{latin-ordinary}{x}}
    <{\OLId{latin-ordinary}{z}} ))]]]$؛
  \item $\lforall[{\OLId{latin-ordinary}{x}}][\lforall[{\OLId{latin-ordinary}{y}}][({\OLId{latin-ordinary}{x}}< {\OLId{latin-ordinary}{y}} \lor \eq[{\OLId{latin-ordinary}{x}}][{\OLId{latin-ordinary}{y}}] \lor {\OLId{latin-ordinary}{y}} < {\OLId{latin-ordinary}{x}})]]$؛
  \item $\lforall[{\OLId{latin-ordinary}{x}}][\lexists[{\OLId{latin-ordinary}{y}}][{\OLId{latin-ordinary}{x}} < {\OLId{latin-ordinary}{y}}]]$؛
  \item $\lforall[{\OLId{latin-ordinary}{x}}][\lexists[{\OLId{latin-ordinary}{y}}][{\OLId{latin-ordinary}{y}} < {\OLId{latin-ordinary}{x}}]]$؛
  \item $\lforall[{\OLId{latin-ordinary}{x}}][\lforall[{\OLId{latin-ordinary}{y}}][({\OLId{latin-ordinary}{x}} < {\OLId{latin-ordinary}{y}} \lif \lexists[{\OLId{latin-ordinary}{z}}][({\OLId{latin-ordinary}{x}} < {\OLId{latin-ordinary}{z}} \land
        {\OLId{latin-ordinary}{z}} < {\OLId{latin-ordinary}{y}}))]]]$.
 \end{enumerate}
\end{defn}

\begin{thm}\ollabel{thm:cantorQ}
  كل ترتيبين خطيين كثيفين !!{enumerable} بلا طرفين متشاكلان.
\end{thm}

\begin{proof}
  لنأخذ $\Struct{M_1}$ و$\Struct{M_2}$ ترتيبين خطيين كثيفين
  !!{enumerable} بلا طرفين، ولنضع ${<_{\OLMathNumeral{1}}} = \Assign{<}{{\OLId{latin-looped}{M}}_{\OLMathNumeral{1}}}$ و
  ${<_{\OLMathNumeral{2}}} = \Assign{<}{{\OLId{latin-looped}{M}}_{\OLMathNumeral{2}}}$. ولتكن $\PIso{I}$ جميع التشاكلات
  الجزئية بينهما. فـ $\PIso{I}$ غير خالية، لأن
  $\emptyset \in \PIso{I}$. نريد إثبات الذهاب والإياب لـ $\PIso{I}$؛
  إذ ينتج منهما $\Struct{M_1} \iso[{\OLId{latin-ordinary}{p}}] \Struct{M_2}$، ثم يكفي
  الرجوع إلى \olref[pis]{thm:p-isom1}.

  أما الذهاب لـ $\PIso{I}$، فنأخذ ${\OLId{latin-ordinary}{p}} \in \PIso{I}$ و
  ${\OLId{latin-ordinary}{a}} \in \Domain{{\OLId{latin-looped}{M}}_{\OLMathNumeral{1}}}$. إن كانت ${\OLId{latin-ordinary}{p}}$ خالية، اخترنا أي
  ${\OLId{latin-ordinary}{b}} \in \Domain{{\OLId{latin-looped}{M}}_{\OLMathNumeral{2}}}$ ووضعنا ${\OLId{latin-ordinary}{q}} = \{\langle {\OLId{latin-ordinary}{a}},{\OLId{latin-ordinary}{b}}\rangle\}$.
  وإن كان ${\OLId{latin-ordinary}{a}}$ في مجال ${\OLId{latin-ordinary}{p}}$، كفانا ${\OLId{latin-ordinary}{q}}={\OLId{latin-ordinary}{p}}$. فبقيت حال ${\OLId{latin-ordinary}{p}}$
  غير الخالية و${\OLId{latin-ordinary}{a}}$ خارج مجالها. نكتب ${\OLId{latin-ordinary}{p}}({\OLId{latin-ordinary}{a}}_{\OLId{latin-ordinary}{i}})={\OLId{latin-ordinary}{b}}_{\OLId{latin-ordinary}{i}}$، حيث ${\OLId{latin-ordinary}{i}}={\OLMathNumeral{1}}$،
  \dots،~${\OLId{latin-ordinary}{n}}$، ونرتب المدخلات بحيث
  ${\OLId{latin-ordinary}{a}}_{\OLMathNumeral{1}} <_{\OLMathNumeral{1}} {\OLId{latin-ordinary}{a}}_{\OLMathNumeral{2}} <_{\OLMathNumeral{1}} \cdots <_{\OLMathNumeral{1}} {\OLId{latin-ordinary}{a}}_{\OLId{latin-ordinary}{n}}$، ولا إخلال بذلك بالعموم.
  ثم نختار ${\OLId{latin-ordinary}{b}} \in \Domain{{\OLId{latin-looped}{M}}_{\OLMathNumeral{2}}}$ بحسب الحالات الآتية:
  \begin{enumerate}
  \item إذا كان ${\OLId{latin-ordinary}{a}} <_{\OLMathNumeral{1}} {\OLId{latin-ordinary}{a}}_{\OLMathNumeral{1}}$، فليكن ${\OLId{latin-ordinary}{b}} \in \Domain{{\OLId{latin-looped}{M}}_{\OLMathNumeral{2}}}$ بحيث
    ${\OLId{latin-ordinary}{b}} <_{\OLMathNumeral{2}} {\OLId{latin-ordinary}{b}}_{\OLMathNumeral{1}}$؛
  \item إذا كان ${\OLId{latin-ordinary}{a}}_{\OLId{latin-ordinary}{n}} <_{\OLMathNumeral{1}} {\OLId{latin-ordinary}{a}}$، فليكن ${\OLId{latin-ordinary}{b}} \in \Domain{{\OLId{latin-looped}{M}}_{\OLMathNumeral{2}}}$ بحيث
    ${\OLId{latin-ordinary}{b}}_{\OLId{latin-ordinary}{n}} <_{\OLMathNumeral{2}} {\OLId{latin-ordinary}{b}}$؛
 \item إذا كان ${\OLId{latin-ordinary}{a}}_{\OLId{latin-ordinary}{i}} <_{\OLMathNumeral{1}} {\OLId{latin-ordinary}{a}} <_{\OLMathNumeral{1}} {\OLId{latin-ordinary}{a}}_{{\OLId{latin-ordinary}{i}}+{\OLMathNumeral{1}}}$ لبعض ${\OLId{latin-ordinary}{i}}$، فليكن
   ${\OLId{latin-ordinary}{b}} \in \Domain{{\OLId{latin-looped}{M}}_{\OLMathNumeral{2}}}$ بحيث ${\OLId{latin-ordinary}{b}}_{\OLId{latin-ordinary}{i}} <_{\OLMathNumeral{2}} {\OLId{latin-ordinary}{b}} <_{\OLMathNumeral{2}} {\OLId{latin-ordinary}{b}}_{{\OLId{latin-ordinary}{i}}+{\OLMathNumeral{1}}}$.
  \end{enumerate}
  ولا يعسر وجود ${\OLId{latin-ordinary}{b}}$ على الوجه المطلوب، فإن $\Struct{M_2}$
  كثيف بلا طرفين. نضع ${\OLId{latin-ordinary}{q}} = {\OLId{latin-ordinary}{p}} \cup \{ \langle {\OLId{latin-ordinary}{a}}, {\OLId{latin-ordinary}{b}} \rangle \}$،
  فتكون ${\OLId{latin-ordinary}{q}} \in \PIso{I}$ وتوسّع ${\OLId{latin-ordinary}{p}}$ كما أردنا. فثبت الذهاب،
  والإياب نظيره. وعليه $\Struct{M_1} \iso[{\OLId{latin-ordinary}{p}}] \Struct{M_2}$؛
  وتنتج من \olref[pis]{thm:p-isom1} العلاقة
  $\Struct{M_1} \iso \Struct{M_2}$.
\end{proof}

\begin{prob}
  أتم برهان \olref[mod][bas][dlo]{thm:cantorQ} بالتحقق من أن
  $\PIso{I}$ تحقق خاصية الإياب.
\end{prob}

\begin{rem}
  كل ترتيب خطي كثيف $\Struct{S}$، إذا كان !!{enumerable} وبلا
  طرفين، لزم من \olref{thm:cantorQ} أن $\Struct{S} \iso \Struct{Q}$،
  حيث $\Struct{Q} = (\Rat, <)$ هو الترتيب الخطي الكثيف
  !!{enumerable} على الأعداد النسبية~$\Rat$. ولنعد إلى
  الـ!!{structure}~$\Struct{R} = (\Real, <)$ في
  \olref[thm]{remark:R}. فقد وجدنا !!a{enumerable}
  !!{structure}~$\Struct{S}$ تحقق $\Struct{R} \elemequiv \Struct{S}$.
  وهذه $\Struct{S}$ ترتيب خطي كثيف !!a{enumerable} بلا طرفين،
  فهي متشاكلة مع الـ!!{structure}~$\Struct{Q}$، ومن ثم مكافئة
  لها أوليًا. وبالتعدي نحصل على $\Struct{R} \elemequiv \Struct{Q}$.
  ولنا طريق مباشر إلى ذلك: نثبت $\Struct{R} \iso[{\OLId{latin-ordinary}{p}}] \Struct{Q}$
  بحجة الذهاب والإياب المتقدمة نفسها.
\end{rem}
\end{document}
